% Image credits for the photographs and AI illustrations of Book 1 (Primary and Middle School).
% Machine-readable license records live in images/book1/CREDITS.md; keep the
% two files in sync. The Book 1 writer owns this file: add one \item per
% photograph (inside the itemize, by chapter) as photographs land.
\chapter*{Image Credits}

The drawings, maps and charts in this book are original; maps and charts
drawn from published data name their source in the caption or in this
list. Photographs are used under their respective free licenses, with
thanks to their authors; full source links and license URLs for every
photograph are listed in \texttt{images/book1/CREDITS.md} in the book's
repository.

\bigskip
Alongside the photographs and the original drawings, some figures are
AI-generated illustrations, produced with OpenAI image generation from
prompts written by the book's editors, and reviewed for scientific
accuracy before inclusion. They show generic scenes only, never a real
place, a real person or a real specimen. The full prompt for every
generated image is recorded in \texttt{images/book1/ai/PROMPTS.md} in the
repository.

\bigskip
\noindent Photographs and data, by chapter:
\begin{itemize}\small
  \item Ch.~\ref{ch:g2:water-on-earth} (water on Earth): the Earth seen from Apollo~17, 7~December 1972 (``The Blue Marble''); NASA, public domain (via Wikimedia Commons). The 100 squares are computed from NOAA ETOPO1 relief; the 100 buckets from the U.S. Geological Survey Water Science School table (after Shiklomanov, 1993).
  \item Ch.~\ref{ch:g3:stones-pebbles-sand} (stones, pebbles and sand): a dry granite stone wall; photograph by Tony Atkin, Geograph, CC BY-SA 2.0 (via Wikimedia Commons). Stone-slab roofs in Switzerland; photograph by Michael C. Rygel, CC BY-SA 3.0 (via Wikimedia Commons).
  \item Ch.~\ref{ch:g3:continents-and-oceans} (the globe): the two globes and the world map are drawn from Natural Earth coastlines (public domain).
  \item Ch.~\ref{ch:g4:clouds-rain-snow} (clouds, rain and snow): cumulus over Kansas (S.~Corfidi), cirrus patches (R.~F.~Kresge), a stratus layer around mountain peaks, a thunderstorm over the Arizona desert (collection of B.~Hooke), a very large hailstone; all NOAA Photo Library, public domain (via Wikimedia Commons).
  \item Ch.~\ref{ch:g4:volcanoes-earthquakes-first-look} (volcanoes and earthquakes): lava entering the ocean at Kilauea, USGS Hawaiian Volcano Observatory; Mount St.~Helens, 18~May 1980, photograph by Austin Post, USGS; a house thrown off its foundation by the 1989 Loma Prieta earthquake, photograph by J.~K.~Nakata, USGS; all public domain (via Wikimedia Commons). The map uses the Smithsonian Global Volcanism Program's list of Holocene volcanoes and the U.S. Geological Survey earthquake catalogue.
  \item Ch.~\ref{ch:g4:world-physical-map} (the physical map): the Himalaya and the Amazon photographed from the International Space Station, NASA, public domain (via Wikimedia Commons). The physical map is drawn from NOAA ETOPO1 relief, with rivers, lakes and coastlines from Natural Earth (public domain).
  \item Ch.~\ref{ch:g5:water-shapes-the-land} (water shapes the land): the Mamoré River photographed from the International Space Station (ISS007-E-10797) and the Nile delta seen by the MODIS instrument on the Terra satellite, NASA, public domain (via Wikimedia Commons).
  \item Ch.~\ref{ch:g5:weather-and-climate} (weather and climate): the Sonoran Desert, U.S. National Park Service, Saguaro National Park, public domain (via Wikimedia Commons). The climate charts are drawn from the NOAA NCEI U.S. Climate Normals 1991--2020 (public domain).
  \item Ch.~\ref{ch:g5:natural-hazards} (natural hazards): the eye of Hurricane Isabel photographed from the International Space Station by Ed Lu, NASA, public domain; the 2014 landslide in Washington State, Mark Reid, U.S. Geological Survey, public domain (both via Wikimedia Commons).
  \item Ch.~\ref{ch:g6:earth-a-planet} (the Earth, a planet): Venus from Mariner~10 (NASA/JPL-Caltech); the Earth from Apollo~17 (NASA); \emph{Earthrise} from Apollo~8 (NASA, Bill Anders); Mars and Jupiter from the Hubble Space Telescope (NASA, ESA, the Hubble Heritage Team and A.~Simon); all public domain (via Wikimedia Commons). The \emph{Pale Blue Dot} seen by Voyager~1 in 1990, 2020 version (NASA/JPL-Caltech, PIA23645, public domain). Planet data from the NASA planetary fact sheet.
  \item Ch.~\ref{ch:g6:latitude-longitude} (latitude and longitude): Harrison's H4 at the National Maritime Museum, Greenwich, photograph by Colonel Warden, CC BY-SA 3.0 (via Wikimedia Commons). Coastlines and cities from Natural Earth (public domain).
  \item Ch.~\ref{ch:g6:rocks-and-minerals} (rocks and minerals): granite, basalt, fossiliferous limestone, sandstone, marble, shale, a quartz crystal and a sheet of muscovite mica, photographs by James St.~John, CC BY 2.0; a slate roof, photograph by {\O}yvind Holmstad, CC0 (all via Wikimedia Commons).
  \item Ch.~\ref{ch:g6:inside-the-earth} (inside the Earth): Inge Lehmann in 1932, photograph by Even Neuhaus, Royal Library of Denmark, public domain (via Wikimedia Commons). The layers, densities and wave speeds are drawn from the Preliminary Reference Earth Model (A.~M. Dziewonski and D.~L. Anderson, 1981).
  \item Ch.~\ref{ch:g6:water-underground} (water underground): Big Spring, Missouri, photograph by Johnston9494, CC BY-SA 4.0 (via Wikimedia Commons). Porosities after R.~C. Heath, \emph{Basic Ground-Water Hydrology}, U.S. Geological Survey (1983); fresh-water shares from the USGS Water Science School.
  \item Ch.~\ref{ch:g7:plate-tectonics} (plate tectonics): the southern continents rejoined, with the fossil bands of \emph{Cynognathus}, \emph{Lystrosaurus}, \emph{Mesosaurus} and \emph{Glossopteris}; U.S. Geological Survey, \emph{This Dynamic Earth}, public domain.
  \item Ch.~\ref{ch:g7:plate-tectonics}: the San Andreas fault across the Carrizo Plain; photograph by Robert E. Wallace, U.S. Geological Survey, public domain.
  \item Ch.~\ref{ch:g7:plate-tectonics}: plate boundaries after P.~Bird (2003), GeoJSON by H.~Ahlenius (ODC-By 1.0); coastlines from Natural Earth (public domain); earthquakes from the U.S. Geological Survey catalogue; volcanoes from the Smithsonian Global Volcanism Program.
  \item Ch.~\ref{ch:g7:volcanoes} (volcanoes): Mount St.~Helens from Johnston Ridge on 17~May and 10~September 1980, photographs by Harry Glicken, U.S. Geological Survey; a lava fountain at Kilauea, 1983, and Mauna Loa from the air, 1985, photographs by J.~D. Griggs, U.S. Geological Survey; the eruption column of Mount Pinatubo, 12~June 1991, photograph by Dave Harlow, U.S. Geological Survey; all public domain (via Wikimedia Commons). The map and the counts use the Smithsonian Global Volcanism Program's list of Holocene volcanoes; plate boundaries after P.~Bird (2003).
  \item Ch.~\ref{ch:g7:earthquakes} (earthquakes): a village near the coast of Sumatra after the tsunami of 26~December 2004, U.S. Navy photograph by Philip A. McDaniel; the fault scarps of the 1983 Borah Peak earthquake, photograph by R.~E. Wallace, U.S. Geological Survey; the road slumped into Hebgen Lake, 1959, photograph by I.~J. Witkind, U.S. Geological Survey; all public domain (via Wikimedia Commons). The counts and depths use the U.S. Geological Survey earthquake catalogue; the intensity table is shortened from the U.S. Geological Survey's Modified Mercalli scale.
  \item Ch.~\ref{ch:g7:sediments-and-fossils} (sediments and fossils): the Grand Canyon from Hermits Rest, NPS photograph by Kristen M. Caldon, public domain; an ammonite, a trilobite and ripple marks in Capitol Reef National Park, photographs by James St.~John, CC BY 2.0; dinosaur footprints at Dinosaur Ridge, photograph by Dexter Perkins, CC0 (all via Wikimedia Commons).
  \item Ch.~\ref{ch:g7:geological-time} (the ages of the Earth): James Hutton, portrait by Henry Raeburn, public domain (via Wikimedia Commons). The time scale's ages and colours are those of the International Chronostratigraphic Chart (International Commission on Stratigraphy).
  \item Ch.~\ref{ch:g8:the-atmosphere} (the atmosphere): launching a weather balloon in Topeka, Kansas, photograph by Stephen Corfidi, NOAA/National Weather Service; the edge of the Earth at sunset from the International Space Station, NASA, ISS Expedition~23; both public domain (via Wikimedia Commons). The profiles are computed from the \emph{U.S. Standard Atmosphere, 1976} (NOAA, NASA and U.S. Air Force).
  \item Ch.~\ref{ch:g8:winds-global-circulation} (winds and the global circulation): the Americas seen by the GOES-16 satellite on 26~July 2022, NOAA, public domain (via Wikimedia Commons). The coastlines of the trade-wind map are from Natural Earth (public domain).
  \item Ch.~\ref{ch:g8:fronts-and-weather-maps} (fronts and weather maps): a depression off Iceland seen by the Terra satellite, NASA/GSFC, MODIS Rapid Response Team, Jacques Descloitres, public domain (via Wikimedia Commons). Coastlines of the air-mass map from Natural Earth (public domain).
  \item Ch.~\ref{ch:g8:the-oceans} (the oceans): HMS \emph{Challenger} preparing to sound, 1872, engraving from the reports of the Challenger expedition reproduced in W.~A. Herdman, \emph{Founders of Oceanography} (1923); the Gulf Stream seen by the MODIS instrument on the Aqua satellite, NASA Earth Observatory; both public domain (via Wikimedia Commons). The ocean-floor profiles are drawn from NOAA ETOPO1 relief; coastlines of the gyre and buoy maps from Natural Earth (public domain).
  \item Ch.~\ref{ch:g8:ice-on-earth} (ice on Earth): Pedersen Glacier in 1990 and 2016, photographs by M.~Tetreau and D.~Kurtz, U.S. National Park Service; the edge of an ice shelf off the Antarctic Peninsula, NASA, Operation IceBridge (M.-J. Vi\~nas); all public domain (via Wikimedia Commons). The sea-ice graph uses the NSIDC Sea Ice Index, version~4 (F.~Fetterer and others, 2025, doi:10.7265/a98x-0f50).
  \item Ch.~\ref{ch:g9:climate-past-and-present} (the Earth's climate, past and present): a section of the GISP2 ice core from Greenland, NOAA Paleoclimatology, public domain (via Wikimedia Commons). The records are drawn from NASA GISTEMP version~4, the Mauna Loa carbon dioxide record (NOAA Global Monitoring Laboratory and Scripps Institution of Oceanography), the EPICA Dome~C temperature (J.~Jouzel and others, 2007), the Antarctic ice-core carbon dioxide composite (B.~Bereiter and others, 2015) and the NOAA Laboratory for Satellite Altimetry sea level; the assessment statements are those of the IPCC Sixth Assessment Report, Working Group~I, Summary for Policymakers (2021).
  \item Ch.~\ref{ch:g9:natural-risks} (living with natural risks): the hazard map of Nevado del Ruiz, from T.~L. Wright and T.~C. Pierson, \emph{Living with Volcanoes}, U.S. Geological Survey Circular 1073 (1992); a tsunami warning buoy inspected by the crew of the USCGC \emph{William Hart}, U.S. Coast Guard, 2021; both public domain (via Wikimedia Commons).
  \item Ch.~\ref{ch:g9:earth-resources} (resources from the Earth): the Chuquicamata mine seen by ASTER, NASA/METI/AIST/Japan Space Systems and the U.S./Japan ASTER Science Team, 2007, public domain; a coal seam in a roadcut in Belmont County, Ohio, photograph by James St.~John, CC BY 2.0 (both via Wikimedia Commons). The concentration factors are from D.~Perkins and others, \emph{Mineralogy}, and S.~Earle, \emph{Environmental Geology}.
  \item Ch.~\ref{ch:g9:landforms} (landforms): the Seven Sisters, photograph by Arild V\aa gen, 2017, CC BY-SA 4.0; crescent dunes on the coast of southern Brazil, NASA astronaut photograph from the International Space Station, 2017, public domain; stalactites in Carlsbad Cavern, National Park Service photograph by Daniel Leifheit, 2018, public domain (all via Wikimedia Commons).
  \item Ch.~\ref{ch:g9:maps-of-the-earth} (maps of the whole Earth): the Lena River delta, Landsat image of 2000, USGS EROS Data Center via NASA Visible Earth, public domain; a GPS satellite, NASA artist's conception, public domain (both via Wikimedia Commons). Coastlines on the world maps from Natural Earth (public domain).
\end{itemize}
